% Generated draft: requires mathematical and source review.
\documentclass[11pt]{article}
\usepackage[margin=25mm]{geometry}
\usepackage{fontspec}
\setmainfont{Latin Modern Roman}
\usepackage{amsmath,amssymb,mathtools,unicode-math}
\setmathfont{Latin Modern Math}
\usepackage{enumitem,needspace,xurl,hyperref}
\setlength{\parindent}{0pt}
\setlength{\parskip}{5pt}
\setlength{\emergencystretch}{3em}
\newcommand{\N}{\mathbb{N}}
\newcommand{\Z}{\mathbb{Z}}
\newcommand{\Q}{\mathbb{Q}}
\newcommand{\R}{\mathbb{R}}
\newcommand{\C}{\mathbb{C}}
\newcommand{\F}{\mathbb{F}}
\newcommand{\A}{\mathbb{A}}
\newcommand{\PP}{\mathbb{P}}
\newcommand{\E}{\mathbb{E}}
\newcommand{\eps}{\varepsilon}
\newcommand{\dd}{\,\mathrm{d}}
\begin{document}
\section*{UnsolvedMath ID: 20002439}
\textbf{Source problem code:} catalogue entry
\subsection{Definitions and mathematical statement}
\textbf{Recovered open problem.} Are nonstationary spherical waves, for example focusing inviscid shocks, stable? The supplied boundary recovery attributes this first-page problem to S.~Benzoni-Gavage, H.~K.~Jenssen and M.~Williams. It specifies no governing gas/system, equation of state, dimension, shock profile, perturbation class or stability norm. The next-presenter kinetic equation is separate material.

\textbf{Supplied self-similar kinematic reduction.} Assume a translation-invariant whole-space problem has a focusing spherical front
\[
 R(t)=A(T-t)^\alpha,\qquad A>0,\quad\alpha>0,\quad t<T,
 \qquad\tau=-\log(T-t)
\]
in normalized time units. A linearized relative radial displacement with spherical harmonic $Y_\ell$ has the form
\[
 \frac{\delta R(t,\vartheta)}{R(t)}
       =\eps e^{s\tau}Y_\ell(\vartheta),\qquad s\in\C.
\]
Its absolute displacement is proportional to $(T-t)^{\alpha-s}$. Therefore relative shape decays for $\operatorname{Re}s<0$, whereas absolute displacement decays for $\operatorname{Re}s<\alpha$. In particular $0<\operatorname{Re}s<\alpha$ is relative growth with absolute shrinking. Equalities are threshold cases; Jordan/logarithmic corrections are not classified. These are linearized magnitude statements, not a nonlinear stability conclusion after a growing relative perturbation ceases to be small.

\textbf{Supplied gauge diagnostics.} Varying the focusing time gives $\delta R/R=\alpha\,\delta T/(T-t)$, hence the radial branch $(\ell,s)=(0,1)$. Moving the center by a constant vector gives first-order absolute displacement $\delta a\cdot\vartheta$, hence the translation branch $(\ell,s)=(1,\alpha)$. These modes must appear in an unprojected similarity formulation \emph{when its admissible spectral domain includes the corresponding parameter derivatives}. Fixing a boundary/center/time or choosing weighted spaces may exclude them, so an arbitrary reduced spectrum is not required to retain them. Stability modulo center and focusing-time changes must modulate or project these branches rather than count them automatically as physical shape instabilities.

The aid reports rigorous ideal-gas Guderley existence results and radial linear/nonspherical modal analyses, but no general nonlinear nonspherical stability theorem in the literature it checked. No specific equation, admissible-parameter theorem or proving citation is supplied. The source problem remains broader than the similarity norm and gauge-mode consistency test.

\textbf{Separate kinetic paragraph.} The appended P.~Szmolyan ``Kinetic and Boltzmann Equations'' setup has damaged typography. Its standard reading would be $f_t+v\cdot\nabla_x f=Q(f,f)$; for a planar profile in direction $e$ the corresponding equation is $(v\cdot e-c)F'=Q(F,F)$. The printed $(v-c)f'$ can be read in one space dimension. The collision operator and kinetic hypotheses are not defined, and no completed kinetic question appears here. This paragraph and the running ARCC/page header are not part of the spherical-shock problem.
\subsection{Short English statement}
The source asks broadly about stability of focusing spherical shocks. The aid supplies a norm-sensitive diagnostic: relative and absolute perturbations have different decay thresholds, and time/center changes produce exact gauge branches. This does not prove nonlinear nonspherical stability. The appended Boltzmann setup belongs to the next presenter.
\subsection{Sources}
\begin{itemize}
\item Local UnsolvedMath registry, permanent ID 20002439, source code catalogue entry; canonical \texttt{aim-pdes-notes.json}, notes[50].
\item AIM/ARCC \emph{Stability criteria for multi-dimensional waves and patterns}, first-page open problem under S.~Benzoni-Gavage, H.~K.~Jenssen and M.~Williams, followed by P.~Szmolyan's separate kinetic section: \url{https://aimath.org/WWN/multidimwaves/multidimwaves.pdf}.
\item R.~Courant and K.~O.~Friedrichs, \emph{Supersonic Flow and Shock Waves}, Interscience Publishers, New York, 1948; L.~D.~Landau and E.~M.~Lifshitz, \emph{Fluid Mechanics}, Pergamon Press, Oxford, 1959, as supplied.
\item No full references are supplied for the later Guderley existence/modal results. Registry research attempt 1, dated 2026-08-14: partial result, status at run open, rigor ``proved'', novelty confidence low. The explicitly machine-generated, non-peer-reviewed aid requires independent verification before citation.
\end{itemize}
\end{document}
